\section{Methodology}\label{sec:methodology}

We trained three pre-trained models on EuroSAT under one shared protocol, tested them on
Aguascalientes without local training, and then fine-tuned them on the Aguascalientes fine-tune
split. This section describes the models, the protocol, the experiments and the metrics.

\subsection{Models and initialization}

We compared ResNet-50, a CNN whose layers learn residual functions;
EfficientNet-B0, the baseline of a CNN family scaled in depth, width
and resolution; and ViT-S/16, a pure transformer applied to a sequence
of 16$\times$16-pixel image patches. This gives two CNN families and one transformer that can
all be trained on one desktop machine with the same settings. We chose the small ViT because its
size is close to that of ResNet-50: with RGB input it has \ParamsViTRgb~M parameters against
\ParamsResNetRgb~M, while EfficientNet-B0 has \ParamsEffNetRgb~M (Table~\ref{tab:cost}). A CNN
versus ViT gap is therefore not just a gap in model size.

All three models started from ImageNet weights (timm). For ResNet-50 and
ViT-S/16 we also used SSL4EO-S12 weights~\cite{wang2023ssl4eo}, learned without labels on
the Sentinel-2 L1C images of the SSL4EO-S12 corpus with self-supervised methods that include MoCo
and DINO. We used the MoCo weights for ResNet-50 and the DINO weights for ViT-S/16. These
weights exist only for the 13 L1C bands, and they expect TorchGeo's normalization, $x/\num{10000}$. EfficientNet-B0 has no
SSL4EO-S12 weights, so it has only ImageNet runs. The SSL4EO-S12 band order puts B8A right after B08, while our data keep B8A
last. We therefore permuted the input channels of the first layer to EuroSAT order; without
this step, B8A and B09 to B12 would each enter through a filter learned for a different band.

ImageNet models expect three input channels, so the 13-band runs needed a new first layer (the
first convolution, or the patch embedding of the ViT). We copied the pre-trained RGB filters to
B04, B03 and B02 and set every other channel to the mean of the three. In this way, the RGB
filters stay exactly as pre-trained, and the new channels start from a neutral filter instead
of random noise that would disturb the pre-trained features in the first epochs.

\subsection{Shared training protocol}

The three models share everything except the architecture, as Table~\ref{tab:config} shows. We
never changed a hyperparameter for one model only, so that differences come from the
architecture and not from tuning. For example, if only the ViT had received a tuned learning
rate, a ViT advantage could be the result of that tuning. Patches were resized from 64 to
224~pixels because this is the input size of the pre-trained weights. ImageNet runs used the
per-band mean and standard deviation of EuroSAT train, also on Aguascalientes, so that no
statistic of the test region enters the pipeline. Every run used seed 0 and was trained once,
so we have no variance estimate across seeds.

\begin{table}[t]
\centering
\footnotesize
\caption{Configuration shared by every run, from the base configuration file. Only the architecture, the input bands and the initial weights change between runs.}
\label{tab:config}
\begin{tabular}{@{}lp{0.68\columnwidth}@{}}
\toprule
Setting & Value \\
\midrule
Input & 64$\times$64 resized to 224$\times$224 (bilinear) \\
Normalization & ImageNet: per-band mean and std of EuroSAT train; SSL4EO-S12: $x/\num{10000}$ (TorchGeo) \\
Augmentation & horizontal flip, vertical flip, 90$^\circ$ rotations \\
Optimizer & AdamW, lr \num{1e-04}, weight decay 0.05 \\
Schedule & 1 warmup epoch + cosine \\
Batch size & 64 (2$\times$32, accumulation) \\
Epochs (EuroSAT) & max 15, patience 3 on val macro-F1 \\
Loss & EuroSAT: cross-entropy; Aguascalientes: class-weighted cross-entropy \\
Fine-tuning & modes \texttt{head}, \texttt{last\_block}; 10 fixed epochs, last epoch kept \\
Seed & 0 \\
\bottomrule
\end{tabular}
\end{table}


We ran every experiment on the same Mac mini with an Apple M2 Pro, using the PyTorch MPS
backend in fp32. Because of memory, the effective batch of 64 was built from two micro-batches
of 32 with gradient accumulation; BatchNorm therefore saw 32 samples at a time.

\subsection{Experiments}

The study has four experiments. In E1 we trained the three models on EuroSAT with ImageNet
weights, each with RGB and with 13 bands (6 runs). In E2 we trained ResNet-50 and ViT-S/16 with
SSL4EO-S12 weights and 13 bands (2 runs). In E3 we evaluated the 8 resulting checkpoints on the
Aguascalientes test split without any local training, that is, zero-shot. In E4 we fine-tuned
the 8 checkpoints on the fine-tune split in two modes (16 runs). Because SSL4EO-S12 runs also
use a different normalization, a gap between SSL4EO-S12 and ImageNet runs cannot be credited to
the weights alone.

The two fine-tuning modes differ in how much of the network they train. In \texttt{head} mode,
the backbone was frozen and only the classifier was trained. In \texttt{last\_block} mode, the
last stage or block was trained together with the classifier. Frozen parts ran in evaluation
mode, so BatchNorm kept its EuroSAT statistics and dropout was off. To keep rare classes from
being ignored, the cross-entropy loss used the class weights
\begin{equation}
  w_c = \frac{n}{K\, n_c},
  \label{eq:class_weights}
\end{equation}
where $n$ is the number of fine-tune patches, $K$ is the number of classes and $n_c$ is the
number of fine-tune patches of class $c$. The weights came from the fine-tune split only.
Fine-tuning ran for a fixed 10 epochs and we kept the last epoch. We chose fixed epochs
because there is no validation split in Aguascalientes: early stopping on the test split would
leak test information into the model. One confound remains: \texttt{last\_block} trains a
different share of the parameters of each architecture, so a CNN versus ViT gap after
\texttt{last\_block} cannot be credited to the architecture alone.

\subsection{Metrics}

Our main metric is macro-F1, the average of the per-class F1 scores. With it, a rare class
counts as much as a common one; with accuracy, a model could score well on Aguascalientes by
predicting HerbaceousVegetation for most patches. The Aguascalientes macro-F1 averages the 9
classes with at least 10 test patches, so River is reported separately. We also report the
macro-F1 over all 10 classes and the accuracy. For cost, we recorded parameters, training time,
peak memory, inference throughput (patches/s) and the epochs needed to reach 95\,\% of the best
validation macro-F1.

To measure the domain shift, we defined the drop and the recovery as
\begin{align}
  \mathrm{Drop} &= \mathrm{F1}_{\mathrm{EuroSAT}} - \mathrm{F1}_{\mathrm{Ags}}^{\text{zero-shot}},
  \label{eq:drop}\\
  \mathrm{Recovery} &= \frac{\mathrm{F1}_{\mathrm{Ags}}^{\text{fine-tuned}}
    - \mathrm{F1}_{\mathrm{Ags}}^{\text{zero-shot}}}{\mathrm{Drop}},
  \label{eq:recovery}
\end{align}
where $\mathrm{F1}_{\mathrm{EuroSAT}}$ is the macro-F1 over the same 9 classes, computed from
the per-class F1 on the EuroSAT test set. In other words, the drop is what a model loses when
it moves to Aguascalientes, and the recovery is the fraction of that loss that fine-tuning wins
back. Recovery is undefined when $\mathrm{Drop} \le 0$.

We fixed the selection rule before seeing any result. The selected model is the one with the
highest Aguascalientes macro-F1 after \texttt{last\_block} fine-tuning. If two models are within
0.01 macro-F1, the one with the higher patches/s on the EuroSAT test set wins. Freezing the
rule first prevents choosing the criterion that favors a preferred model after the fact.

\subsection{Grad-CAM}

To see which parts of a patch drive a prediction, we used Grad-CAM~\cite{selvaraju2017}. We
computed it for the predicted class on the 224$\times$224 input, resized the map to the
64$\times$64 patch and scaled it to $[0, 1]$. The target layer was the last block of
\texttt{layer4} for ResNet-50 and \texttt{conv\_head} for EfficientNet-B0, the last
convolutional features of each CNN. For ViT-S/16 we used the first LayerNorm (\texttt{norm1})
of the last block. There, we dropped the CLS token and reshaped the 196 patch tokens into a
14$\times$14 map, so the gradients have a spatial layout. We chose this layer because it is the
last point where the patch tokens still affect the prediction: after the last attention layer,
only the CLS token reaches the classifier, so the gradients on the patch tokens would be zero.
